\begin{statementbox}
	\textbf{\textcolor{CStatement}{条件}}
	\begin{description}[style=nextline, leftmargin=2.8em, labelsep=0.5em, itemsep=0.35em]
		\item[\textbf{\textcolor{CStatement}{条件1（二次函数预设）：}}]
			$f(x)=ax^2+bx+c$（$a\neq0$），记 $\Delta=b^2-4ac$，对称轴 $x_0=-\frac{b}{2a}$，开区间 $(m,n)$。
	\end{description}
	\textbf{\textcolor{CStatement}{结论}}
	\begin{description}[style=nextline, leftmargin=2.8em, labelsep=0.5em, itemsep=0.45em]
		\item[\textbf{\textcolor{CStatement}{结论一（两不等实根）：}}]
			\textbf{\textcolor{CStatement}{直观描述：}}
			抛物线在开区间内与$x$轴有两个不同交点。
			\[
				\Delta>0,\quad m<x_0<n,\quad af(m)>0,\quad af(n)>0.
			\]
		\item[\textbf{\textcolor{CStatement}{结论二（唯一实根）：}}]
			\textbf{\textcolor{CStatement}{直观描述：}}
			区间两端函数值异号，曲线穿过$x$轴一次。
			\[
				f(m)f(n)<0.
			\]
		\item[\textbf{\textcolor{CStatement}{推广结论（两实根含重根）：}}]
			\textbf{\textcolor{CStatement}{直观描述：}}
			允许对称轴为切点，即两实根可相等。
			\[
				\Delta\ge0,\quad m<x_0<n,\quad af(m)>0,\quad af(n)>0.
			\]
	\end{description}
	\textbf{\textcolor{CStatement}{使用提醒：}}
	推广结论包含重根情形（$\Delta=0$），适用于题目不要求“不等根”的情况。原结论一要求$\Delta>0$，仅适用不等根。
\end{statementbox}
